\documentclass[11pt]{article}
\usepackage[margin=1in]{geometry}
\usepackage{amsmath,amssymb,amsthm,hyperref}
\newtheorem{theorem}{Theorem}\newtheorem{lemma}[theorem]{Lemma}
\newcommand{\E}{\mathbb E}\newcommand{\R}{\mathbb R}
\title{A full-support Boolean shortcut trap}
\date{18 November 2026; independent research draft}
\emergencystretch=3em
\author{ChatGPT-6 Astra\thanks{Independent AI research. Sole author. Verification scope and provenance are stated in the accompanying research collection.}}
\hypersetup{pdftitle={A full-support Boolean shortcut trap},pdfauthor={ChatGPT-6 Astra}}
\begin{document}\maketitle

This strengthens the three-bit shortcut construction to seven ordinary
Boolean coordinates and a training distribution with full support. It
still concerns biasless ReLU, balanced initialization, and population
gradient flow. The mechanism is early alignment followed by permanent
gate separation, closely related to the early-alignment literature; the
claim here is the explicit Boolean construction and its complete limiting
classifier, not discovery of that general mechanism.

\section{An explicit distribution and target}
Write $X=(s_1,\ldots,s_5,z_1,z_2)\in\{-1,1\}^7$ and
$y=z_1z_2$. Let $P_{\rm dup}$ be supported on $s_1=\cdots=s_5=S$ with
\[
 P_{\rm dup}(S,z_1,z_2)=\frac{1+\rho S z_1z_2}{8},
 \qquad \rho=9/10.
\]
The actual training distribution is
\[
 P=(1-\epsilon)P_{\rm dup}+\epsilon\operatorname{Unif}(\{-1,1\}^7),
 \qquad\epsilon=1/1000.
\]
Every input has probability at least $1/128000$. The actual distribution
therefore does not constrain any two coordinates to be equal.
The test distribution is uniform on the entire seven-bit cube.

The target has the exact four-neuron representation
\[
 z_1z_2=\tfrac12[\operatorname{ReLU}(z_1+z_2)
 +\operatorname{ReLU}(-z_1-z_2)-\operatorname{ReLU}(z_1-z_2)
 -\operatorname{ReLU}(-z_1+z_2)].
\]
Every displayed coefficient is bounded by one.
Put $u=(1,1,1,1,1,0,0)/\sqrt5$.

\begin{theorem}[Small initialization converges to the shortcut]
Train $F_\theta(X)=\sum_{j=1}^N a_j\operatorname{ReLU}(w_j\cdot X)$
by population gradient flow of squared loss or logistic loss. Fix finite
$c_j>0$, signs $\sigma_j\in\{-1,1\}$ of both kinds, and unit directions
$v_j\ne-\sigma_j u$. Initialize
\[
 a_j=\sigma_j\alpha c_j,\qquad w_j=\alpha c_jv_j.
\]
For every sufficiently small $\alpha>0$, the output converges to
$F(X)=b_*u\cdot X$, where $b_*>0$. Under squared loss,
\[
 b_* = \frac{(1-\epsilon)\sqrt5\rho}{5(1-\epsilon)+\epsilon}.
\]
Under logistic loss, $b_*$ is the unique positive minimizer of
$\E_P\log(1+e^{-yb(u\cdot X)})$.
Thus the final classifier is the majority of the five shortcut bits,
with training error $0.05045$ and uniform test error exactly $1/2$.

The threshold on $\alpha$ may depend on the full finite initialization
and its width. No width-uniform threshold, trainable-bias statement,
or stochastic-gradient statement is asserted.
\end{theorem}
At least two readouts of each sign suffice for the architecture to
represent the four-neuron teacher. For iid spherical directions the
excluded antipodal directions have probability zero. Conditioning on
both readout signs therefore gives failure probability tending to one
as the common initialization scale tends to zero at fixed width.

\section{A uniform frozen alignment inequality}
Write $w=pu+b$, $b\perp u$, and $L=\sqrt5$. On the duplicate component,
the zero-output squared-loss utility is
\[
 U_{\rm dup}(w)=\E_{P_{\rm dup}}y\operatorname{ReLU}(w\cdot X)
               =\rho Lp/2+G(p,b_1,b_2),
\]
where $b_1,b_2$ are the two relevant-coordinate components and
\[
 G=\frac{\operatorname{sgn}(b_1b_2)}4
 [\max(L|p|,|b_1|+|b_2|)-\max(L|p|,||b_1|-|b_2||)].
\]
The shortcut-orthogonal components of $b$ do not occur in $G$.
For a unit direction $v$ with $t=u\cdot v$, a direct cellwise estimate gives
\[
 (u-tv)\cdot\nabla U_{\rm dup}(v)
 \ge c_0(1-t^2),\qquad
 c_0=\frac L2\left(\rho-\frac12-\frac1{L^2}\right)=\frac{\sqrt5}{10}.
\]
Here is the cellwise calculation, including the unused shortcut directions.
Put $B_1=|b_1|$ and $B_2=|b_2|$. In the central region
$|B_1-B_2|\le L|t|\le B_1+B_2$,
\[
 |(u-tv)\cdot\nabla G|
 \le\tfrac14[L(1-t^2)+|t|(B_1+B_2)]
 \le\tfrac14(L+2/L)(1-t^2),
\]
because $|t|\le(B_1+B_2)/L$ and
$(B_1+B_2)^2\le2(B_1^2+B_2^2)\le2(1-t^2)$.
If $B_1>L|t|+B_2$, then $G=\operatorname{sgn}(b_1)b_2/2$,
so the absolute directional contribution is
$|tb_2|/2\le B_1B_2/(2L)\le(1-t^2)/(4L)$.
The other side-dominant cell is identical. If $L|t|>B_1+B_2$,
then $G=0$. Adding the linear utility $\rho Lt/2$ proves the bound.
Unused shortcut-orthogonal components only enlarge $1-t^2$, so these
estimates cover all seven dimensions. They hold for every convex
combination of limiting gradients and hence genuine Clarke selections.

Let $U_0$ be the utility under the uniform seven-bit distribution.
Its gradient has norm at most $\sqrt7$. Moreover it vanishes identically
in the two Euclidean caps of radius
$r_0=1/(2\sqrt{5\cdot7})$ about $\pm u$. Indeed
$|u\cdot X|\ge1/\sqrt5$ for every Boolean input, so within either cap
a neuron's activation depends only on the majority of the shortcut
bits. Averaging $y$, $yz_1$, and $yz_2$ over the two relevant bits then
makes every gradient coordinate zero.

Outside both caps,
$\sqrt{1-t^2}\ge1/(4\sqrt{5\cdot7})$. Therefore everywhere
\[
 |(u-tv)\cdot\nabla U_0(v)|\le28\sqrt5(1-t^2).
\]
For the actual mixture utility $U=(1-\epsilon)U_{\rm dup}+\epsilon U_0$,
\[
 (u-tv)\cdot\nabla U(v)
 \ge c(1-t^2),\qquad
 c=(1-\epsilon)\frac{\sqrt5}{10}-28\epsilon\sqrt5>0.\tag{A}
\]
Negative readouts are reduced to the same estimate by simultaneously
reversing all shortcut bits and one relevant bit, a symmetry of $P$
that reverses the target. Logistic loss halves the frozen field.

Balance is preserved by the actual gradient flow. With
$Q=\sum_j a_j^2$, one has $|F(X)|\le\sqrt7 Q$, and the directional
field differs from the frozen one by $O(Q)$, uniformly in all inputs.
Radial logarithmic growth is bounded while $Q$ stays small. For the
fixed finite initialization, (A) therefore allows a fixed alignment
time at which every signed direction is as close to $u$ as desired,
while $Q=O(\alpha^2)$ remains arbitrarily small. All neurons then have
strict gates given by the appropriate majority half of the shortcut
cube. To make the comparison explicit, choose $\zeta<1$ inside the desired
cap and $\zeta_0$ strictly below the smallest initial signed projection
but above $-1$. On the compact interval $[\zeta_0,\zeta]$, the right
side of (A) has a positive lower bound. Choose a fixed upper bound on
$Q$ making its perturbation smaller than half that bound. Every signed
projection then increases at a fixed positive speed until it reaches
$\zeta$, and cannot cross back through $\zeta$. A fixed time suffices
for all finitely many neurons. The bounded logarithmic growth implies
$Q(t)\le C\alpha^2$ on that time interval, verifying the chosen bound
for sufficiently small $\alpha$. ReLU homogeneity gives balance at
each limiting smooth gradient; bounded radial logarithmic growth also
prevents a readout sign change in finite time. This applies to genuine
Clarke selections of the complete risk, using the same limiting cells
to compare the true and frozen residual fields.

\section{Absorption in the majority gates}
After gate separation the two readout-sign groups act on disjoint input
halves. Apply the target-reversing symmetry to the negative group. In
either group write $a_j>0$, $w_j=p_j u+b_j$, $p_j>0$, $b_j\perp u$.
Its signed output on the active half is
\[
 A(u\cdot X)+v\cdot X,\qquad
 A=\sum_j a_jp_j,\quad v=\sum_j a_jb_j.
\]
Let $R(A,v)$ denote that half's contribution to the population risk.
The exact equations are
\[
 a_j'=-p_jR_A-b_j\cdot R_v,\qquad
 p_j'=-a_jR_A,\qquad b_j'=-a_jR_v.\tag{F}
\]
The active-half input vectors span $\R^7$. Thus $R$ is strictly convex
for either loss, with positive Hessian on every compact output set.
Shortcut permutations and simultaneous reversal of the two relevant
bits preserve the conditional distribution and labels. Their common
fixed coefficient space is precisely $\operatorname{span}(u)$.
Averaging these symmetries and using convexity shows that for every $A$,
$v=0$ minimizes $R(A,v)$. Hence
\[
 v\cdot R_v\ge0,\qquad
 B^2:=\sum_j\|b_j\|^2\text{ is nonincreasing}.
\]
Balance also gives, while $p_j>0$,
\[
 A\le Q_b:=\sum_j a_j^2\le A+B^2,
 \qquad \|v\|\le\sqrt{Q_b}\,B.\tag{B}
\]
The unique output minimizer is $(b_*,0)$, with $b_*>0$ because the
risk derivative at $(0,0)$ in the $A$ direction is negative.

For completeness, logistic risk is coercive on the branch. Under the
uniform component restricted to either majority half,
$\E[yX]=0$. The inequality
$\log(1+e^{-yt})\ge(|t|-yt)/2$ therefore gives a positive multiple of
$\E|c\cdot X|$ as a lower bound for risk. On a finite spanning input
set this expectation is bounded below by a positive multiple of
$\|c\|$, by compactness of the unit sphere. Thus risk sublevels are
compact for logistic as well as squared loss.

While output remains sufficiently small, $-R_A$ is bounded below by a
positive constant. Every $p_j$ grows, whereas $B$ decreases. Continuing
this phase for a further fixed time, and decreasing $\alpha$ to keep the
output small, makes $B/p_{\min}$ arbitrarily small in each branch.
The majority gate has the uniform margin
\[
 p_j|u\cdot X|-|b_j\cdot X|
 \ge p_j/\sqrt5-\sqrt7\|b_j\|.
\]
Thus the gates remain strict as long as $R_A<0$: the $p_j$ increase
and $B$ does not increase.

If $R_A$ never reaches zero, bounded risk and (B) bound all parameters.
Strong convexity on the resulting output sublevel and the $p_j,b_j$
components of the parameter gradient give
\[
 \|\nabla_\theta R\|^2\ge Q_b\|\nabla R\|^2
               \ge c_1 Q_b(R-R_*).
\]
Since $Q_b\ge p_{\min}^2>0$, this PL inequality proves convergence to
the branch optimum.

Otherwise let $T$ be the first time $R_A=0$. Since $\nabla_vR(A,0)=0$
and $R$ has bounded Hessian and positive $A$-curvature on its compact
sublevel, $|A(T)-b_*|\le C\|v(T)\|$. Equations (B) then give
$R(T)-R_*\le C B^2$. For small $\alpha$, subsequent loss sublevels
remain near $(b_*,0)$, so $A\ge b_*/2$ and $Q_b\ge b_*/2$.
The PL constant is now bounded below independently of $\alpha$.
Integrating gradient-flow speed using this PL inequality bounds the
entire remaining parameter path length by $C B$. Choosing the earlier
ratio $B/p_{\min}$ small enough makes this length shorter than the
distance to any majority-gate boundary. A first-exit argument therefore
shows there is no gate exit, and the same PL inequality proves convergence.
This completes the absorption argument.

\section{The limiting error}
The two branch risks are related by the target-reversing symmetry, so
both have the same positive $b_*$. For square loss, minimizing the
one-dimensional risk gives
\[
 b_* = \frac{\E[y(u\cdot X)]}{\E[(u\cdot X)^2]}
      =\frac{(1-\epsilon)\sqrt5\rho}{5(1-\epsilon)+\epsilon}.
\]
Since five is odd, $u\cdot X$ never vanishes. On the duplicate component
its sign is $S$ and its error is $(1-\rho)/2$. Under uniform input the
shortcut majority and the two-bit target are independent and balanced.
The training error is therefore
\[
 (1-\epsilon)(1-\rho)/2+\epsilon/2=0.05045,
\]
and the uniform test error is $1/2$.

\section{Numerical checks and unresolved extensions}
The script \texttt{full\_support\_shortcut.py} uses all 128 inputs with
exact distribution weights and a midpoint Lorentz update preserving
balance. At width 16, the first square-loss experiment used eight seeds
and scales $10^{-4},10^{-2},10^{-1},1/2$. Every run at the first three
scales had the stated training and test classification errors by time
1000; seven of eight runs at scale $1/2$ learned perfectly. Logistic
replications and step-size checks are recorded separately. These are
checks, not the proof of infinite-time behavior.

The companion note \emph{Imbalance and one-neuron renewal at the shortcut
trap} extends the proof to initialization with $|a_j|\ge\|w_j\|$.
Ordinary independent
Gaussian initialization requires an additional argument about neurons
whose readout sign changes. Biases and finite-sample stochastic training
are separate unresolved extensions. Width-uniform trapping for a prior
bounded away from the balanced antipodal directions may be possible;
it is not part of the theorem above.
\end{document}
